\section{Discussion et limites}
\label{sec:discussion}

\subsection{Synthèse des résultats}

Les résultats de ce projet confirment, sur une tâche de classification
fine-grained à petite échelle, le comportement data-hungry du ViT
documenté dans la littérature~\cite{dosovitskiy2020vit,touvron2021deit} :
sans pré-entraînement à grande échelle, le ViT n'apprend pas de
représentation exploitable dans les conditions expérimentales de ce
projet, quelle que soit la quantité de données d'entraînement fournie.
À l'inverse, ResNet-50 exploite efficacement son biais inductif
convolutif pour apprendre une représentation utile même avec 10\,\% du
jeu d'entraînement, et atteint 71{,}9\,\% de top-1 avec le jeu complet.
Le ViT pré-entraîné, quant à lui, confirme l'intérêt du pré-entraînement
à grande échelle : sa variante patch32 atteint 53{,}9\,\% de top-1,
là où le ViT from scratch reste bloqué au niveau du hasard.

L'application de démonstration développée pour ce projet illustre ce
contraste de manière concrète : sur une image test correctement
identifiée, ViT et ResNet-50 s'accordent avec des niveaux de confiance
cohérents ; sur une image hors distribution (espèce absente de
CUB-200-2011), les deux modèles produisent des confiances faibles et
sont en désaccord, un comportement souhaitable plutôt qu'une erreur
silencieuse (Section~\ref{sec:deployment}).

\subsection{Limites}

\textbf{Nombre d'epochs restreint.} Les modèles ont été entraînés sur
seulement 3 epochs (Section~\ref{sec:methodology}), une contrainte de
temps de calcul disponible sur la durée du projet. Cette limite affecte
en premier lieu le ViT from scratch, dont l'absence totale
d'apprentissage (accuracy proche du hasard indépendamment de la
quantité de données) est probablement davantage due à un
sous-entraînement qu'à une limite intrinsèque de l'architecture. Un
entraînement plus long, combiné à un scheduler de taux d'apprentissage
et à un arrêt anticipé (tous deux implémentés dans l'outil
d'entraînement du projet mais non encore appliqués aux résultats
rapportés), constituerait la première piste d'amélioration.

\textbf{Écart inexpliqué entre les variantes patch16 et patch32 du ViT
pré-entraîné.} La variante patch16 (9{,}9\,\% de top-1) performe
nettement moins bien que la variante patch32 (53{,}9\,\%), un résultat
contre-intuitif qui n'a pas pu être expliqué dans le temps imparti et
mériterait une investigation plus poussée (vérification du taux
d'apprentissage, du nombre d'itérations, ou d'un éventuel problème de
convergence spécifique à cette configuration).

\textbf{Échelle réduite du ViT from scratch.} L'implémentation
personnalisée du ViT (6 blocs, dimension d'embedding 384) reste
nettement plus petite que le ViT-Base original~\cite{dosovitskiy2020vit},
ce qui limite la portée des comparaisons avec les résultats publiés
dans la littérature.

\textbf{Absence de variance statistique.} Chaque configuration n'a été
entraînée qu'une seule fois (un seul seed), sans répétition pour
estimer la variance des résultats d'un run à l'autre.

\textbf{Taille du jeu de données.} Avec 5\,094 images d'entraînement,
CUB-200-2011 reste un jeu de données modeste au regard des standards
de pré-entraînement des Vision Transformers, ce qui limite la capacité
du ViT from scratch à démontrer son potentiel dans ce cadre expérimental.